\documentclass[fleqn,10pt]{src/KampSjabloon} 

\definecolor{color1}{RGB}{63, 81, 181} % Kleur titels
% Dit kan ook bv. 241,134,34 (oranje), 143,0,0 (rood) zijn

\definecolor{color2}{RGB}{70, 85, 190} % achtergrondkleur
% Dit kan ook bv. 177,114,0 (oranje), 70,0,0 (rood) zijn

%--------------------------------------------------------------------
%  KAMPVARIABELEN
%--------------------------------------------------------------------

% Onderstaande variabelen kunnen doorheen het kamp gebruikt worden
% door bv. \kampprijs tussen de tekst te zetten, deze moet later
% dan slechts op één plek aangepast worden. Onderstaande gegevens 
% worden daarbij ook gebruikt om het financieel overzicht te maken.

\newcommand{\kampprijs}{100} % Bijdrage per lid

\newcommand{\specialeformule}{\fpeval{round((175 + (\aantalnachten-1)*50)/(\ledenaantal*\aantalnachten),2)}} % berekent huurprijs p.p. per nacht volgens de formule van jou lokaal (indien ze niet per persoon rekenen). Geef dan \specialeformule in als huurprijs.

\newcommand{\huurprijs}{6} % Huurprijs per persoon per nacht
\newcommand{\transportprijs}{273} % Transportkost per persoon
\newcommand{\voedselprijs}{10} % Prijs p.p. per dag eten & drank
\newcommand{\kasoverschot}{6500} % Overschot takkas
\newcommand{\materiaalprijs}{30} % Budget spelmateriaal
\newcommand{\kampaandenken}{200} % Budget kampaandenken
\newcommand{\activiteitprijs}{20} % Activiteitprijs p.p. daguitstap,...
\newcommand{\afvalkost}{30} % Geschatte totale vuilzakken kost

\newcommand{\ledenaantal}{18}
\newcommand{\leidingaantal}{4}
\newcommand{\fourageaantal}{0} % Geen transportkost, wel huurprijs

\newcommand{\aantalnachten}{7} % Aantal nachten kamp
\newcommand{\kampplaats}{Timboektoe}

\KampTitel{Jinkamp \the\year} % Kamp titel
\Tak{Jins } % Sjamfoeters , ... (gevolgd door spatie)

\InhetKort{
Hier is plaats voorzien voor een korte uiteenzetting van het kamp, een beschrijving van het kampthema, de uitwerking tijdens het kamp en misschien iets over de setting/locatie van het kamp ...
}

\affiliation{dd/07/2015 uu:mm - dd/07/2015 uu:mm} % Kamptermijn

%--------------------------------------------------------------------
%  Begin Document
%--------------------------------------------------------------------

\begin{document}
\flushbottom 
\maketitle % Plaatst de titel, tak, ...
\tableofcontents % Maakt inhoudstafel aan
\parskip = 1pc
\parindent = 0pc

%--------------------------------------------------------------------
%  KAMPLOCATIE
%--------------------------------------------------------------------

\section{Kamplocatie} % Dit is een hoofdtitel

\begin{figure}[h]\centering
\includegraphics[width=8.5cm]{Images/bulgaria-air.jpg}
\caption{Kampindicatie}
\end{figure}

\subsection*{Bestemming} % Ondertitel, met ster niet in inhoudstafel

Een korte beschrijving van \kampplaats, inleiding tot de bestemming/streek is altijd leuk .. 

\subsection*{Kampterrein}

Uitleg over het soort kampterrein, misschien iets dat leuk of speciaal is aan dit kampterrein of deze kamplocatie. \\ % Deze schuine strepen dwingen enter af.

\point\ Gemeentepark 14, 2930 \kampplaats \\

%--KAMPINFO-------------------------
\begin{table}[H]
\centering
\begin{tabular}{ll}
\toprule
\multicolumn{2}{c}{\textbf{\kampplaats}} \\
\toprule

Officiële landstaal & Bulgaars\\
Hoofdstad & Sofia\\
Regeringsvorm & Parlementaire republiek\\
Religie & Bulgaars-orthodox (85\%)\\
Oppervlakte & 110.879 km2\\
Inwoners & 6.981.642 (2013)\\
Munteenheid & Lev (1 lev = 0,51 euro) \\
UTC	& +2 (zomertijd +3) \\

\bottomrule
\end{tabular}
\label{tab:label}
\end{table}
%-----------------------------------

%--------------------------------------------------------------------
%  ALGEMENE DAGINDELING
%--------------------------------------------------------------------

\section{Algemene dagindeling}

%--DAGINDELING----------------------
\begin{table}[H]
\centering
\begin{tabularx}{5cm}{lX}
\toprule

10h00 & Opstaan\\
11h00 & Voormiddagactiviteit\\
13h00 & Middageten\\
14h30 & Namiddagactiviteit\\
16h00 & Vieruurtje\\
19h00 & Avondeten\\
20h00 & Avondactiviteit\\
22h00 & Tienuurtje\\
00h00 & Gaan slapen\\

\bottomrule
\end{tabularx}
\caption{Algemene dagindeling}
\end{table}
%-----------------------------------

%--------------------------------------------------------------------
%  LEDEN
%--------------------------------------------------------------------

\section{Leden}

We gaan op kamp met \textbf{\ledenaantal} leden, waarvan de meerderheid / minderheid \male / \female. We passen bij de \Tak wel / geen indeling in patrouilles of nestjes  toe. De kostprijs van het kamp per lid bedraagt \textbf{\euro{ \kampprijs}} te betalen vóór dd/mm. Wij bieden verminderde kampprijs aan, aan gezinnen met financiële moeilijkheden.

%--------------------------------------------------------------------
%  LEIDING EN FOURAGE
%--------------------------------------------------------------------

\section{Leiding en Fourage}

%--CONTACTGEGEVENS------------------
\begin{table}[H]
\centering
\begin{tabularx}{8.5cm}{Xccc}
\toprule
\multicolumn{4}{c}{\textbf{Leiding}} \\
\midrule
Naam & GSM & RB?& KV? \\
\midrule

Jonas Verheyen & +32 498420388 & Neen & Ja \\
Rafaël Hautekiet & +32 478339404 & Ja & Neen \\

\midrule
\multicolumn{4}{c}{\textbf{Fourage}} \\
\midrule
Naam & GSM & RB? & TV? \\
\midrule

Sam Verschooren & +32 400000001 & Neen & Neen \\
Sam Verschooren & +32 400000002 & Neen & Neen \\

\bottomrule
\end{tabularx}
\caption{Overzicht leiding en fourage}
\end{table}
%-----------------------------------

%--------------------------------------------------------------------
%  AFSPRAKEN
%--------------------------------------------------------------------

\section{Afspraken}

\subsection{Afspraken met de leden}

%--AFSPRAKEN------------------------
\begin{itemize}

\item We dragen zorg voor elkaar.
\item ...

\end{itemize}
%-----------------------------------

\textit{De \Tak dienen deze regels na te leven, indien dit niet gebeurt zullen er eventuele sancties volgen.} % dit zet de tekst schuin

\subsection{Afspraken met de leiding/fouriers}

%--AFSPRAKEN------------------------
\begin{itemize}

\item No limits
\item ...

\end{itemize}
%-----------------------------------

\subsection*{BOB-regeling}

Details bobregeling, wie wanneer bob is staat in meer detail in het kampprogramma.

%--------------------------------------------------------------------
%  VAN EN NAAR HET KAMPTERREIN
%--------------------------------------------------------------------

\section{Van en naar het kampterrein}
Detail transport van en naar kampterrein/-locatie

%--REISINFO-------------------------
\begin{table}[H]
\centering
\begin{tabularx}{8.5cm}{llXXl}
\toprule
\textbf{Datum} & \multicolumn{2}{c}{\textbf{Vertrek}} & \multicolumn{2}{c}{\textbf{Aankomst}} \\
\toprule

 dd/07/2015 & uu:mm &  Jobertus & \kampplaats & uu:mm \\
 dd/07/2015 & uu:mm &  \kampplaats & Jobertus & uu:mm \\

\toprule
Reisduur & \multicolumn{2}{c}{X uur Y minuten}& & \\
\bottomrule
\end{tabularx}
\caption{Details kampvervoer}
\end{table}
%-----------------------------------

%--------------------------------------------------------------------
%  COMMUNICATIE MET DE OUDERS
%--------------------------------------------------------------------

\section{Communicatie met de ouders}

Details ouderavond, noodnummer, ...

%--------------------------------------------------------------------
%  BELANGERIJKE INFO
%--------------------------------------------------------------------

\section{Belangrijke info}

%--INFO-----------------------------
\begin{table}[H]
\centering
\begin{tabularx}{8.5cm}{lX}
\toprule
\multicolumn{2}{c}{\textbf{Algemene noodnummers}} \\
\toprule

Noodgevallen & 112\\
Antigifcentrum & 070 24 52 45\\
SGV nood & +32474 26 14 01\\

\bottomrule
\multicolumn{2}{c}{\textbf{Noodnummers \kampplaats}} \\
\toprule

Dokter  & Dr. Assen Aleksiev\\
        & 15 Dimitar Nestorov Blvd \\
        & Sofia  Bulgaria \\
        & + 359 2 581 2833\\

Tandarts  & Dr. Todor Peev \\
          & Department of Stomatology \\
          & 1 Georgi Sofiyski Blvd \\
          & Sofia Bulgaria \\
          & + 359 2 851 8832 \\

Ziekenhuis & Tokuda Hospital\\
		   &51B Nikola Vaptzarov Blvd. \\
           &Sofia 1407 Bulgaria \\
           &+ 359 2 943 2215\\

\bottomrule
\end{tabularx}
\label{tab:label}
\end{table}
%-----------------------------------

\subsection*{EHBO}

Wie is verantwoordelijke, wie volgde EHBO cursus? Wat we allemaal meenemen van medicatie zit in bijlage. 

\subsection{Verzekering}

\import{/}{Bijlage/verzekering} % Uitleg verzekering van SGV staat in verzekering.tex
 
%--------------------------------------------------------------------
%  FINANCIEEL OVERZICHT
%--------------------------------------------------------------------

\section{Financieën}

We vertrekken op kamp met een bankkaart van Record Bank met standaad afhaallimiet van \euro{ 600} per dag met een limiet van \euro{ 1200} per week. 

\financieeloverzicht % dit commando maakt overzicht met de kampvariabelen

%--------------------------------------------------------------------
%  UITWERKING KAMPTHEMA
%--------------------------------------------------------------------

\section{UItwerking Kampthema}

Info en uitwerking van je kampthema ...

%--------------------------------------------------------------------
%  HIKE/DAGTOCHT
%--------------------------------------------------------------------

\section{Hike/dagtocht}

Info en uitwerking Hike/dagtocht ...

%--------------------------------------------------------------------
%  TOTEMISATIE/BELOFTEN
%--------------------------------------------------------------------

\section{Totemisatie/Beloften}

Alle info over de totemisatie/beloften, uitwerking in thema, opdrachtjes, ...

\textbf{Opdrachten stad:}

%-----------------------------------
\begin{itemize}

\item Leer jullie zelf voorstellen in het Bulgaars
\item ...

\end{itemize}
%-----------------------------------

%--------------------------------------------------------------------
%  EVALUATIEMOMENT
%--------------------------------------------------------------------

\section{Evaluatiemoment}

Info omtrent evaluatiemoment, hoe pakken jullie dit origineel aan?

\textbf{Voorbereidde stellingen/onderwerpen:}

%-----------------------------------
\begin{itemize}

\item Ik vind dat het geld inzamelen ons jaar teveel overheerst heeft en zou liever minder hebben gewerkt.
\item Ik had graag meer scouteske avontuurlijke activiteiten gedaan tijdens het jaar
\item Ik had graag meer vergaderingen gehad tijdens het jaar, in plaats van sportpaleizen, feestjes, quiz, ...
\item Ik ben tevreden over de groepssfeer.
\item Ik had er gewoon meer van verwacht ... :(

\end{itemize}
%-----------------------------------

%--------------------------------------------------------------------
%  KAMPPROGRAMMA
%--------------------------------------------------------------------

\section{Kampprogramma}

Eventuele inleiding tot kampprogramma ... Onderstaande dagbeschrijvingen zijn vooral om de rode draad in het kamp duidelijk te maken.

%--OVERZICHT------------------------
\begin{table}[h]
\centering
\begin{tabularx}{8.5cm}{lXXX}
\toprule
\textbf{Datum} & \textbf{Voormiddag} & \textbf{Namiddag} & \textbf{Avond}\\
\toprule

\hyperref[11]{11 juli} & Estafette & Vrij spel & Sing-song \twonotes\\
\hyperref[12]{12 juli} & Heenreis & Receptie & Klassespel \\
\hyperref[13]{13 juli} & Knutselen & Klassespel & Casino \\

\toprule
\end{tabularx}
\caption{Overzicht Kampprogramma}
\end{table}
%-----------------------------------

%--DAGBESCHRIJVING------------------
\subsection*{11 juli}
\label{11}
\point\ Sunny Beach \\

\begin{figure}[H]\centering
\includegraphics[width=8.5cm]{Images/sunny.jpg}
\caption{Sunny Beach}
\end{figure}

Beschrijving dagplanning\footnote{Een eventuele voetnoot}, de speluitwerkingen vind je telkens onder 'Speluitwerkingen', een hyperlink naar een bepaalde speluitwerking maak je zo \hyperref[stdsspl]{\it{Stadsspel}}% belangerijk hierbij is dat je \label{stdsspl} onder de titel van de desbetreffende speluitwerking plaatst. 
 of zo \hyperref[runandhide]{\it{Run and Hide}}. Als plan B kan er steeds verwezen worden naar het \hyperref[backup]{Back-up programma}. % Opnieuw kan je dit doen odmat er \label{backup} onder de titel Back-up programma staat

\textcolor{color1}{\checked} Macramégerief, Voetbal, Tekengerief,  ...
%-----------------------------------

%--DAGBESCHRIJVING------------------
\subsection*{12 juli}
\label{12}
\point\ Locatie \\

\begin{figure}[H]\centering
\includegraphics[width=8.5cm]{Images/boat-cruise.jpg}
\end{figure}

Beschrijving dagplanning

\textcolor{color1}{\checked} Zaken die we deze dag nodig hebben ...
%-----------------------------------

%--DAGBESCHRIJVING------------------
\subsection*{13 juli}
\label{13}
\point\ Locatie \\

Beschrijving dagplanning

\textcolor{color1}{\checked} Zaken die we deze dag nodig hebben ...

%--UITGAVENTABEL VOOR TIJDENS REIS--
\begin{table}[H]
\begin{tabular}{lllll}
\toprule

 \textbf{Uitgaven Rila gebergte} & \textbf{Schatting} & \textbf{Kost} \\
\toprule
Trein Jakoruda $\rightarrow$ Kostandovo & \euro{ 50,00} &  ....... \\
3 x Eten ontbijt, middag, avond & \euro{ 500,00} &  .......  \\
\toprule
\textbf{Balans} & \multicolumn{2}{c}{............................}& & \\

\bottomrule
\end{tabular}
\label{tab:label}
\end{table}
%-----------------------------------

%-----------------------------------


%--------------------------------------------------------------------
%  SPELUITWERKINGEN
%--------------------------------------------------------------------

\section{Speluitwerkingen}

%--SPELUITWERKING------------------
\subsection*{Stadsspel}
\label{stdsspl} % Dit label dient om in het kampprogramma te refereren.
% bijvoorbeeld door in de tekst \hyperref[stdsspl]{stadsspel} te plaatsen.

\subsubsection*{Inleiding}
\subsubsection*{Spelverloop}
\subsubsection*{Benodigdheden}

\begin{itemize}
\item Goed humeur
\item ...
\end{itemize}
%-----------------------------------

%--SPELUITWERKING------------------
\subsection*{Run and Hide!}
\label{runandhide} % Dit label dient om in het kampprogramma te refereren.
% bijvoorbeeld door in de tekst \hyperref[stdsspl]{stadsspel} te plaatsen.

\subsubsection*{Inleiding}
\subsubsection*{Spelverloop}
\subsubsection*{Benodigdheden}

\begin{itemize}
\item ...
\end{itemize}
%-----------------------------------

%--------------------------------------------------------------------
%  BACK-UP PROGRAMMA
%--------------------------------------------------------------------

\section{Back-up programma}
\label{backup}
Het kan natuurlijk niet altijd meezitten (regen, storm, ...) en vroeg of laat zou de verveling wel eens kunnen toeslaan, wij zijn vastberaden met dit back-up programma het nooit zover te laten komen.

\import{/}{Bijlage/back-up.tex} % BACK-UP spelletjes staan uitgeschreven in back-up.tex

%--------------------------------------------------------------------
%  BIJLAGE
%--------------------------------------------------------------------

\section{Bijlage}

\subsection{Benodigdheden
}

\subsubsection*{Leiding}

%-----------------------------------
\begin{itemize}

\item Walkie-Talkies
\item EHBO-koffer + medische fiches
\item ...

\end{itemize}
%-----------------------------------

\subsubsection*{Leden}
Info omtrent bagage, afspraken soort zak, buddy systeem, ...?

%-----------------------------------
\begin{itemize}

\item	Kleine rugzakje
\item	Slaapzak
\item	Matje, luchtmatras wordt afgeraden
\item	K-way
\item	Stapschoenen
\item	Sandalen/teensleffers
\item	T-shirts
\item	Broeken (minstens één lange) / rokjes
\item	Trui
\item	Ondergoed
\item	Iets tegen de zon (pet, hoed, ..)
\item	Zwemgerief
\item	Bestek
\item	Gamel
\item	Keukenhanddoek
\item	Ecover afwasmiddel + shampoo 
\item	Stick om kleren te wassen
\item	Tandpasta
\item	Shampoo + zeep
\item	Muggenmelk
\item	Zonnecrème + After-sun
\item	Pillamp
\item	Scoutsuniform (hemd en das)!

\end{itemize}
%-----------------------------------

\subsection{Ledenlijst}

%--LEDENLIJST-----------------------
\begin{table}[H]
\begin{tabularx}{8.5cm}{lXccc}
\toprule  % je kan \square vervangen door \blacksquare om het in te kleuren
\textbf{Voornaam}& \textbf{Achternaam}&& \textbf{Tel. Ouders} & \textbf{\checked}\\
\toprule

Toon      & Degenaers    &\male& 0400000000 & $\square$ \\ \midrule
Laura     & Buschgens    &\female& 0400000000 & $\square$ \\ \midrule
Jarne     & De Kort      &\male& 0400000000 & $\blacksquare$ \\ \midrule
Kaat      & Welkenhuysen &\female& 0400000000 & $\square$ \\ \midrule
Joke      & Daems        &\female& 0400000000 & $\square$ \\ \midrule
Caroline  & Bernaerts    &\female& 0400000000 & $\square$ \\ \midrule
Loes      & Cardoen      &\female& 0400000000 & $\square$ \\ \bottomrule

\end{tabularx}
\caption{Ledenlijst \Tak}
\end{table}
%-----------------------------------

\newpage % hiermee zorg je dat hij op een nieuwe pagina begint, indien gewenst.

\subsection{EHBO}
\import{/}{Bijlage/ehbo} % EHBO staat uitgeschreven in ehbo.tex

% voorbeeld van pagina vullende afbeelding
\begin{figure*}[h]
\centering
\includegraphics[width=17.5cm]{Images/voorbeeld.png}
\caption{Stadspel Sofia: Team Groen Krantenartikel 4}
\end{figure*}

% voorbeeld van afbeelding in landschap stand
\begin{landscape}
\centering
\begin{figure}[h]
\centering
\vspace{1.5cm}
\includegraphics[width=23cm]{Images/railmap.pdf}
\caption{Spoorwegnet Bulgarije}
\end{figure}
\end{landscape}

\end{document}